\documentclass[12pt,letterpaper,oneside]{report}

% Kompilator w Overleaf: pdfLaTeX
\usepackage[utf8]{inputenc}
\usepackage[T1]{fontenc}
\usepackage[polish]{babel}
\usepackage{polski}
\usepackage[letterpaper,left=2.0cm,right=2.0cm,top=2.0cm,bottom=2.0cm]{geometry}
\usepackage{microtype}
\setlength{\parindent}{1.25em}
\setlength{\parskip}{0pt}
\setlength{\emergencystretch}{2em}

\usepackage{amsmath,amssymb,mathtools,bm}
\usepackage{booktabs}
\usepackage{graphicx}
\graphicspath{{figures/}}
\usepackage{caption}
\captionsetup{font=small,labelfont=normalfont,justification=justified,singlelinecheck=false}
\usepackage{float}
\usepackage{titlesec}
\titleformat{\chapter}[display]
  {\normalfont\Large\bfseries}{\chaptername\ \thechapter}{0.85em}{}
\titlespacing*{\chapter}{0pt}{0pt}{2.2em}
\titleformat{\section}{\normalfont\large\bfseries}{\thesection}{0.9em}{}
\titlespacing*{\section}{0pt}{1.7em}{0.65em}
\renewcommand{\thesection}{\arabic{section}}
\numberwithin{equation}{chapter}
\numberwithin{figure}{chapter}
\numberwithin{table}{chapter}

\usepackage{tocloft}
\renewcommand{\contentsname}{Spis treści}
\renewcommand{\cfttoctitlefont}{\hfill\Large\bfseries}
\renewcommand{\cftaftertoctitle}{\hfill}
\setlength{\cftbeforetoctitleskip}{0.5cm}
\setlength{\cftaftertoctitleskip}{1.0cm}

\newcommand{\UniversityName}{Uniwersytet Jagielloński w Krakowie}
\newcommand{\FacultyName}{Wydział Fizyki, Astronomii i Informatyki Stosowanej}
\newcommand{\AuthorName}{Karina Nepravskaya}
\newcommand{\AlbumNumber}{1209730}
\newcommand{\SupervisorName}{tytuł, imię i nazwisko promotora}
\newcommand{\DepartmentName}{nazwa zakładu lub instytutu}
\newcommand{\YearOfSubmission}{2026}
\newcommand{\Pe}{\ensuremath{P_e}}
\newcommand{\Neff}{\ensuremath{N_{\mathrm{eff}}}}
\newcommand{\xmax}{\ensuremath{x_{\max}}}
\newcommand{\direct}{\textit{direct}}
\newcommand{\delayed}{\textit{delayed}}

\usepackage{hyperref}
\hypersetup{colorlinks=true,linkcolor=blue,citecolor=blue,urlcolor=blue,
  pdftitle={Analiza korelacji pędów elektronów w potrójnej jonizacji neonu},
  pdfauthor={\AuthorName}}

\pagestyle{plain}
\setcounter{secnumdepth}{1}
\setcounter{tocdepth}{1}

\begin{document}

\begin{titlepage}
\thispagestyle{empty}
\noindent{\bfseries \UniversityName}\par
\noindent\FacultyName

\vspace{2.7cm}
\begin{center}
\Large \AuthorName

\vspace{0.35cm}
{\large Nr albumu: \AlbumNumber}
\end{center}

\vspace{2.0cm}
\begin{center}
{\LARGE\bfseries Analiza korelacji pędów elektronów\\[0.2em]
w potrójnej jonizacji neonu\\[0.2em]
w silnym polu laserowym\par}
\end{center}

\vfill
\begin{flushright}
\begin{minipage}{0.48\textwidth}
\raggedleft
Praca wykonana pod kierunkiem\\
\SupervisorName\\
\DepartmentName
\end{minipage}
\end{flushright}

\vfill
\begin{center}Kraków \YearOfSubmission\end{center}
\end{titlepage}

\pagenumbering{arabic}
\setcounter{page}{2}
\tableofcontents
\clearpage

\phantomsection
\addcontentsline{toc}{chapter}{Streszczenie}
\begin{center}\textbf{Streszczenie}\end{center}
\vspace{1.2em}

Praca dotyczy korelacji końcowych pędów trzech elektronów emitowanych w
potrójnej jonizacji neonu w silnym, liniowo spolaryzowanym polu laserowym.
Punktem wyjścia jest analiza głównych składowych pełnego wektora pędów.
Wskazuje ona jeden podłużny kierunek ruchu wspólnego oraz dwa kierunki
opisujące różnice pomiędzy elektronami. Ponieważ osie PCA zależą od
konkretnej próbki, do właściwej analizy wprowadzono stałą, ortonormalną
transformację Jacobiego--Helmerta. Rozdziela ona całkowity pęd elektronów
od ich względnej nierównowagi, nie tracąc informacji zawartej w trzech
składowych podłużnych.

W tej reprezentacji porównano kanały \direct{} i \delayed{} dla natężenia
\(1{,}7\,\mathrm{PW/cm^2}\). Kanał \direct{} cechuje się silniejszymi
dodatnimi korelacjami pędów, mniejszą nierównowagą względną i bardziej
równomiernym podziałem wartości bezwzględnych pędu. Wynik ten jest zgodny
z bardziej kolektywnym charakterem tego kanału, jednak same pędy końcowe
nie stanowią samodzielnego dowodu przebiegu czasowego jonizacji.

\chapter{Wstęp}

W potrójnej jonizacji nie wystarcza pytanie o to, jaki pęd osiąga
pojedynczy elektron. Równie istotne jest, jak w obrębie jednego zdarzenia
dzieli się pęd pomiędzy trzy elektrony. Taka informacja jest rozproszona
w trzech składowych i łatwo ją zgubić, gdy każdą z nich analizuje się
osobno. Celem pracy jest zatem skonstruowanie opisu, który rozdziela ruch
wspólny całego układu elektronowego od różnic pomiędzy jego częściami, a
następnie wykorzystanie tego opisu do porównania dwóch kanałów jonizacji.

Pole laserowe jest spolaryzowane wzdłuż osi \(z\), dlatego właśnie
podłużne składowe pędu zawierają zasadniczą informację o ruchu dryfowym.
W dalszym tekście zapisano je skrótowo jako
\begin{equation}
p_2\equiv p_{2z},\qquad p_3\equiv p_{3z},\qquad p_4\equiv p_{4z}.
\label{eq:longitudinal-momenta}
\end{equation}
W obliczeniach trajektoryjnych elektrony 2 i 3 są oznaczone jako
początkowo związane, a elektron 4 jako tunelujący. Są to jednak etykiety
historii obliczenia, nie zaś fizycznie rozróżnialne własności elektronów.
Z tego powodu wielkości zależne od numeracji są używane tylko jako opis
trajektorii, a główne wnioski opierają się również na miarach symetrycznych
względem permutacji elektronów.

Wcześniejsza reprezentacja simpleksowa opisuje względne udziały
\begin{equation}
x_i=\frac{|p_i|}{|p_2|+|p_3|+|p_4|},\qquad x_2+x_3+x_4=1.
\label{eq:shares}
\end{equation}
Dla trzech elektronów każdy zestaw udziałów można przedstawić jako punkt
w trójkącie. Jest to użyteczne, ale wykres nie zawiera znaków pędów ani
ich skali. Ponadto dla większej liczby elektronów trójkąt zastępuje
simpleks wyższego wymiaru, który nie daje już równie czytelnej
wizualizacji. Potrzebny jest więc opis, w którym różne aspekty korelacji
nie są sztucznie upychane w jednym obrazie.

\chapter{Opis korelacji pędowych}

Analizę rozpoczęto od PCA dziewięciu końcowych składowych pędu trzech
elektronów. W danych wyróżnia się kierunek z podobnymi współczynnikami
przy \(p_{2z}\), \(p_{3z}\) i \(p_{4z}\), odpowiadający ich wspólnemu
ruchowi wzdłuż osi polaryzacji. Dwa pozostałe niezależne kierunki
podłużne opisują różnice pomiędzy pędami elektronów. To rozpoznanie jest
ważne, ale PCA nie może być końcową definicją osi: wektory własne zależą
od kowariancji konkretnego zbioru, a dwa kierunki w tej samej podprzestrzeni
mogą zmienić orientację przy zmianie kanału lub próbki.

Z tego powodu dla każdego zdarzenia zastosowano stałą bazę
Jacobiego--Helmerta:
\begin{equation}
Q=\frac{p_2+p_3+p_4}{\sqrt{3}},\qquad
\xi_1=\frac{p_2-p_3}{\sqrt{2}},\qquad
\xi_2=\frac{p_2+p_3-2p_4}{\sqrt{6}}.
\label{eq:helmert}
\end{equation}
Współrzędna \(Q\) opisuje ruch kolektywny. Współrzędne \(\xi_1\) i
\(\xi_2\) tworzą płaszczyznę ruchu względnego: pierwsza porównuje
elektrony 2 i 3, druga porównuje tę parę z elektronem 4. Transformacja
nie redukuje wymiaru, lecz jedynie zmienia bazę, o czym świadczy równość
\begin{equation}
p_2^2+p_3^2+p_4^2=Q^2+\xi_1^2+\xi_2^2.
\label{eq:norm}
\end{equation}
Można więc porównywać kanały w tym samym układzie współrzędnych, zachowując
pełną informację o podłużnych pędach.

Wielkość
\begin{equation}
K=\sqrt{\xi_1^2+\xi_2^2}
\label{eq:imbalance}
\end{equation}
jest miarą względnej nierównowagi. Jest równa zeru tylko wtedy, gdy trzy
pędy są jednakowe, a jej wartość nie zależy od numeracji elektronów. Samo
\(K\) nie określa jednak kierunku nierównowagi; tę informację zachowuje
para \((\xi_1,\xi_2)\). Do opisu podziału wartości bezwzględnych
wykorzystano ponadto efektywną liczbę uczestniczących elektronów
\begin{equation}
\Neff=\frac{1}{x_2^2+x_3^2+x_4^2}
\label{eq:neff}
\end{equation}
oraz największy udział \(\xmax=\max(x_2,x_3,x_4)\). Duże \(\Neff\) i
małe \(\xmax\) oznaczają równomierniejszy podział. Znaki pędów są
analizowane osobno w ośmiu sektorach, na przykład \texttt{+++} lub
\texttt{++-}; dzięki temu nie ginie informacja o kierunku emisji.

\chapter{Porównanie kanałów jonizacji}

Porównanie wykonano dla tej samej intensywności pola,
\(I=1{,}7\,\mathrm{PW/cm^2}\), co pozwala przypisać różnice przede
wszystkim klasyfikacji trajektorii, a nie zmianie warunków zewnętrznych.
W tabeli~\ref{tab:comparison} zestawiono miary, które opisują odmienne
aspekty zdarzenia: szerokość ruchu kolektywnego, nierównowagę względną,
podział wartości bezwzględnych oraz korelację pary elektronów 2 i 3.

\begin{table}[htbp]
\centering
\caption{Porównanie kanałów dla \(I=1{,}7\,\mathrm{PW/cm^2}\).}
\label{tab:comparison}
\begin{tabular}{lrrrrrr}
\toprule
kanał & \(n\) & \(\sigma(\Pe)\) & \(\langle K\rangle\) &
\(\langle\Neff\rangle\) & \(\langle\xmax\rangle\) & \(r_{23}\)\\
\midrule
\direct  & 1695 & 7,543 & 1,398 & 2,628 & 0,473 & 0,868\\
\delayed & 1601 & 5,384 & 2,383 & 2,374 & 0,533 & 0,214\\
\bottomrule
\end{tabular}
\end{table}

Kanał \direct{} ma mniejszą średnią wartość \(K\), większe
\(\Neff\) i mniejsze \(\xmax\). Te trzy obserwable opisują ten sam
trend z różnych stron: w kanale \direct{} pęd jest dzielony bardziej
równomiernie, a trzy elektrony pozostają bliżej wspólnego kierunku ruchu.
Silna dodatnia korelacja \(r_{23}\) wzmacnia tę interpretację. Dla kanału
\delayed{} rozkład w przestrzeni względnej jest szerszy, korelacja pary
jest słabsza, a pojedynczy elektron częściej przejmuje większą część
bezwzględnego pędu.

Ten sam kontrast pojawia się w sektorach znakowych. W kanale \direct{}
sektor \texttt{+++} obejmuje \(47{,}4\%\) zdarzeń, podczas gdy w kanale
\delayed{} jego udział wynosi \(23{,}0\%\). Przewaga emisji o zgodnych
znakach jest więc spójna z bardziej kolektywnym ruchem pierwszego kanału.
Nie należy jednak utożsamiać tej zgodności z bezpośrednim dowodem
mechanizmu czasowego: końcowy pęd jest wynikiem całej trajektorii i sam
nie zastępuje analizy chwil jonizacji.

\chapter{Podsumowanie}

PCA ujawniła w danych strukturę złożoną z jednego kierunku kolektywnego
i dwóch kierunków względnych. Jej rolą było rozpoznanie tej struktury,
a nie wyznaczenie ostatecznych osi analizy. Stała transformacja
Jacobiego--Helmerta pozwala zapisać tę samą informację w układzie, który
ma prostą interpretację i pozostaje wspólny dla porównywanych zbiorów.

Wyniki dla \(1{,}7\,\mathrm{PW/cm^2}\) konsekwentnie odróżniają kanały
\direct{} i \delayed{}. Pierwszy z nich ma bardziej kolektywną strukturę
pędów i równomierniejszy ich podział, drugi zaś większe wewnętrzne
zróżnicowanie. Zaproponowany opis nie zastępuje pełnego rozkładu
wielociałowego, lecz porządkuje go w zestaw niewielkiej liczby
uzupełniających się obserwabli. Dzięki temu stanowi wygodny punkt wyjścia
do dalszego badania zależności od intensywności pola oraz do rozszerzenia
analizy na większą liczbę elektronów.

\end{document}
